\documentclass[]{article}

\usepackage{amsmath}
\usepackage{multirow}
\usepackage{natbib}
\usepackage{graphicx} 


\begin{document}

\subsection{Numerical simulation dataset}The analysis is based on a dataset from a high-resolution direct numerical simulation of three-dimensional, homogeneous, and isothermal turbulence. The dataset consists of 1001 snapshots of the velocity field, equally spaced in time. The simulation domain is a periodic cube, and the physical spacing between grid points is $dx=0.0078125$. Prior to analysis, each velocity field snapshot was processed with a Gaussian low-pass filter ($\sigma=1.0$ voxel) to mitigate numerical noise at the grid scale. The velocity gradient tensor, $\nabla \mathbf{v}$, was computed from the filtered velocity fields using a second-order finite difference scheme, with periodic boundary conditions handled appropriately.\subsection{Vortex identification and tracking}Coherent vortex structures were identified in each snapshot using the Q-criterion, which isolates regions of the flow where the rate-of-rotation tensor, $\Omega_{ij} = \frac{1}{2}(\partial_j v_i - \partial_i v_j)$, dominates the rate-of-strain tensor, $S_{ij} = \frac{1}{2}(\partial_j v_i + \partial_i v_j)$  \citep{menon2021significancestraindominatedregionvortex}. The criterion is defined as:\begin{equation}Q = \frac{1}{2} \left( ||\Omega||^2 - ||S||^2 \right) > 0\end{equation}where $||\cdot||$ denotes the Frobenius norm. To ensure the robustness of our findings, we identified vortex structures using a range of thresholds for $Q$, defined relative to the mean ($\mu_Q$) and standard deviation ($\sigma_Q$) of the Q-field in each snapshot, specifically $Q > \mu_Q + k\sigma_Q$ with $k \in \{2.5, 3.0, 4.0\}$. Only connected regions with a volume of at least 20 voxels were considered as distinct vortex structures.For each identified vortex, we computed its vorticity-weighted centroid, $\mathbf{r}_c$, to define its position  \citep{kashyap2026emergencelocalorderinggiant}. These centroids were then tracked across consecutive snapshots to construct Lagrangian trajectories. The tracking was performed using a greedy nearest-neighbor algorithm, linking a centroid at time $t$ to the closest centroid at time $t+1$ within a search radius of three grid cells, accounting for the periodic domain using the minimum-image convention. Only trajectories persisting for a minimum of 20 snapshots were retained for the final statistical analysis to ensure sufficient data for characterizing their long-time behavior.\subsection{Statistical analysis of trajectories}To characterize the transport regime, we performed a multi-faceted statistical analysis on the ensemble of vortex trajectories.First, we computed the time-averaged Mean Squared Displacement (MSD) for a time lag $\tau$  \citep{bailey2022logitfitactive}:\begin{equation}\text{MSD}(\tau) = \langle |\mathbf{r}(t+\tau) - \mathbf{r}(t)|^2 \rangle  \citep{katayama2025notetwopointmeansquare}\end{equation}where $\mathbf{r}(t)$ is the position of a vortex centroid at time $t$ and the angle brackets denote an average over all possible start times $t$ and all trajectories. The diffusive exponent, $\alpha$, was determined by fitting a power law, $\text{MSD}(\tau) \propto \tau^\alpha$, to the data in a log-log plot. This analysis was performed for the total displacement vector as well as for its individual Cartesian components to probe for anisotropy in the diffusion process  \citep{mcclure2026collectivevortexdynamicsisolated}.Second, to investigate the temporal correlations in the vortex motion, we calculated the normalized Velocity Autocorrelation Function (VACF):  \citep{bradley2025velocitycorrelationsvorticesrarefaction}\begin{equation}C_v(\tau) = \frac{\langle \mathbf{v}(t) \cdot \mathbf{v}(t+\tau) \rangle}{\langle |\mathbf{v}(t)|^2 \rangle}\end{equation}where the vortex velocity $\mathbf{v}(t)$ was estimated using a finite difference of the centroid positions  \citep{ashraf2025experimentalinvestigationwaterexit}. The decay rate of the VACF provides a measure of the memory in the vortex's motion.Third, we analyzed the statistics of the trajectory step sizes to test for deviations from a simple random walk. The step displacement vectors, $\Delta \mathbf{r}_i = \mathbf{r}_{i+1} - \mathbf{r}_i$, were computed for all trajectories. The probability distribution of the step magnitudes, $|\Delta \mathbf{r}|$, was then compared against a Gaussian distribution. The goodness-of-fit was quantified using the Kolmogorov-Smirnov (K-S) test.\subsection{Evaluation of transport anisotropy and coupling}To connect the statistical properties of the transport to the underlying physics, we performed two additional analyses. First, we quantified the coupling between a vortex's motion and the surrounding fluid. The background fluid velocity was interpolated at the location of each vortex centroid. We then computed the Pearson correlation coefficient between the components of the vortex velocity and the corresponding components of the local fluid velocity.Second, we investigated the anisotropy of vortex motion relative to its own orientation. The local vorticity vector, $\boldsymbol{\omega} = \nabla \times \mathbf{v}$, was computed at each vortex centroid. For each step displacement $\Delta\mathbf{r}$, we decomposed it into components parallel ($\Delta\mathbf{r}_\parallel$) and perpendicular ($\Delta\mathbf{r}_\perp$) to the local vorticity vector $\boldsymbol{\omega}$. The mean magnitudes of these components were then compared to determine if there is a preferential direction of motion relative to the vortex filament's axis  \citep{zhang2019imaginganisotropicvortexdynamics}.

\end{document}
                